% !PRINTED 18
\weberpage{18}
%<<CONT
ausfallen. Es gilt aber der folgende schöne und wichtige Satz
von C.\,\gesperrt{Jordan} $^{1)}$:

\gesperrt{I. Wie auch die Compositionsreihe einer Gruppe
$P$ gewählt sein mag, die Indexreihe ist, von
der Anordnung abgesehen, immer dieselbe.}

Der Beweis beruht auf der Betrachtung der im vorigen
Paragraphen eingeführten complementären Gruppen.

Es seien $Q$ und $Q_1$ Normaltheiler von $P$ und zugleich $Q_1$
ein Theiler (also nach \S.\,4, 5. Normaltheiler) von $Q$.

Dann gilt der Satz:

\gesperrt{1. Die Gruppe $Q/Q_1$ ist ein Normaltheiler der
Gruppe $P/Q_1$, und der Index von $Q/Q_1$ in Bezug
auf $P/Q_1$ ist gleich dem Index von $Q$ in Bezug
auf $P$.}

Um seine Richtigkeit einzusehen, braucht man nur die
Bildungsweise der einzelnen Gruppen genauer zu betrachten.

Es seien $n, q, q_1$ die Grade von $P, Q, Q_1$ und

\weq{3}{n = \nu q,\quad q = \nu_1 q_1,\quad n = \mu q_1,\quad \mu = \nu \nu_1.}

Man zerlegt nun $Q$ in die Nebengruppen zu $Q_1$, d. h. man
setzt, wenn $b_1, b_2 \ldots b_{\nu_1 - 1}$ passend bestimmte Elemente in $Q$ sind,

\weq{4}{Q = Q_1 + Q_1 b_1 + \cdots Q_1 b_{\nu_1 - 1}.}

Die Elemente der Gruppe $Q/Q_1$ sind

\weq{5}{Q_1,\, Q_1 b_1,\, \ldots\, Q_1 b_{\nu_1 - 1}.}

Wenn wir $P$ in die Nebengruppen zu $Q_1$ zerlegen, so kommen
 darunter sicher die Elemente (5) vor, und folglich ist $Q/Q_1$
ein Theiler von $P/Q_1$.

Es ist also noch nachzuweisen, dass dieser Theiler ein
Normaltheiler ist.

Bezeichnen wir aber mit $a$ irgend ein Element von $P$, so
sind alle Elemente von $P/Q_1$ in der Form $Q_1 a$ enthalten, und
wir haben also nur zu zeigen, dass, wenn $b$ irgend ein Element
in $Q$ ist, jedes Element

\weq{6}{Q_1 a^{-1} Q_1 b Q_1 a}

in $Q/Q_1$, d. h. in der Form $Q_1 b$ enthalten ist.

$^{1)}$ C.\,\gesperrt{Jordan}, Traité des substitutions, Paris 1870. \gesperrt{Netto}, Substitutionentheorie,
Leipzig 1882. Wir folgen hier dem Beweise, den
\gesperrt{Hölder} im Anschluss an \gesperrt{Netto} gegeben hat (Math. Annalen, Bd. 34).
